% Part: second-order-logic
% Chapter: sol-and-set-theory
% Section: power-of-continuum

\documentclass[../../../include/open-logic-section]{subfiles}

\begin{document}

\olfileid{sol}{set}{pow}

\olsection{સાતત્યકનું કદ}

\begin{explain}
દ્વિતીય-ક્રમ તર્કશાસ્ત્રમાં !!{domain}ના ઉપગણો પર પરિમાણન કરી
શકીએ છીએ, પરંતુ !!{domain}ના ઉપગણોના ગણો પર સીધું પરિમાણન કરી
શકતા નથી. તે માટે \emph{તૃતીય-ક્રમ} તર્કશાસ્ત્ર જોઈએ. દાખલા તરીકે,
ગણના શક્તિસમૂહમાંથી તે ગણમાં કોઈ !!{injective} વિધેય નથી એવો
કૅન્ટૉરનો પ્રમેય જણાવવો હોય, તો આપણે ``દરેક ગણ~$X$ અને દરેક
ગણ~$P$ માટે, જો~$P$ એ~$X$નો શક્તિસમૂહ હોય, તો
$\cardle{P}{X}$ નહીં'' એમ લખવાનો પ્રયાસ કરી શકીએ. પરંતુ~$P$
એ~$X$નો શક્તિસમૂહ છે એમ કહેવા માટે~$P$ના !!{element}s એ~$X$ના
બધા અને માત્ર ઉપગણો છે એમ ઔપચારિક રીતે વ્યક્ત કરવું પડે; દાખલા
તરીકે $\lforall[Y][(P(Y) \liff Y \subseteq X)]$ જેવું કંઈક.
અડચણ~$P(Y)$માં છે: તે દ્વિતીય-ક્રમ તર્કશાસ્ત્રનું !!a{formula}
નથી, કારણ કે~$P$ જેવા એકસ્થાની સંબંધ-ચલોના દલીલો માત્ર પદો હોઈ
શકે છે.

છતાં, વ્યાપકક્ષેત્ર પૂરતું મોટું હોય તો ગણોના ગણો પરના પરિમાણનનું
\emph{અનુકરણ} કરી શકીએ છીએ. વિચાર એ છે કે દ્વિસ્થાની સંબંધ~$R$
!!{domain}ના !!{element}sને !!{domain}ના જ !!{element}s સાથે જોડે છે.
એવો~$R$ હોય તો કોઈ~$x$ સાથે~$R$-સંબંધિત બધા !!{element}sનો ગણ લઈએ:
$\Setabs{y \in \Domain{M}}{R(x,y)}$ એ~$x$ દ્વારા
``સંકેતિત'' ગણ છે. ઊલટું, $Z \subseteq \Pow{\Domain{M}}$
એ~$\Domain{M}$ના ઉપગણોનો કોઈ ગણ હોય અને~$\Domain{M}$માં
$Z$ના ગણોની સંખ્યા જેટલા અથવા વધુ !!{element}s હોય, તો એવો
સંબંધ~$R \subseteq \Domain{M}^2$ પણ મળે છે કે દરેક $Y \in Z$
કોઈ~$x$ દ્વારા~$R$નો ઉપયોગ કરીને સંકેતિત થાય.
\end{explain}

\begin{defn}
જો $R \subseteq \Domain{M}^2$ હોય, તો \emph{$x$ એ~$R$ વડે}
$\Setabs{y \in \Domain{M}}{R(x, y)}$ને \emph{સંકેતિત કરે છે}.
\end{defn}

જો !!a{element}~$x \in \Domain{M}$ એ~$R$ વડે
$Z \subseteq \Domain{M}$ ગણને સંકેતિત કરે, તો
$Y \subseteq \Domain{M}$ ગણોના ગણને, એટલે~$Y$ના !!{element}s
દ્વારા સંકેતિત ગણોને, સંકેતિત કરે છે. તેથી ગણ~$Y$ એ~$R$ વડે
$\Pow{X}$ને સંકેતિત કરી શકે છે. એવું ત્યારે અને માત્ર ત્યારે જ બને
જ્યારે દરેક $Z \subseteq X$ માટે કોઈ $x \in Y$ એ~$R$ વડે~$Z$ને
સંકેતિત કરે અને દરેક $x \in Y$ એ~$R$ વડે કોઈ $Z \subseteq X$ને
સંકેતિત કરે.

\begin{prop}
!!{formula}
  \[
  \fn{Codes}(x, R, Z) \ident \lforall[y][(Z(y) \liff
    R(x, y))]
  \]
એ વ્યક્ત કરે છે કે $s(x)$ એ~$s(R)$ વડે~$s(Z)$ને સંકેતિત કરે છે.
નીચેનું !!{formula}
\begin{multline*}
  \fn{Pow}(Y, R, X) \ident \\
  \lforall[Z][(Z \subseteq X \lif \lexists[x][(Y(x) \land
    \fn{Codes}(x, R, Z))])] \land {} \\ \lforall[x][(Y(x) \lif
  \lforall[Z][(\fn{Codes}(x, R, Z) \lif Z \subseteq X)])]
\end{multline*}
એ વ્યક્ત કરે છે કે $s(Y)$ એ~$s(R)$ વડે~$s(X)$ના શક્તિસમૂહને
સંકેતિત કરે છે: $s(Y)$ના ઘટકો~$s(R)$ વડે~$s(X)$ના બધા અને માત્ર
ઉપગણોને સંકેતિત કરે છે.
\end{prop}
\sourcecorrection{OLSOL-018}{સ્થિર અંગ્રેજી મૂળના~$\fn{Pow}$
સૂત્રના બીજા સંયોજકમાં~$Y(x)$થી શરૂ થતો કૌંસ બંધ થતો નથી.
અહીં અંતે એક બંધ કૌંસ ઉમેર્યો છે, જેથી બહારનું નિહિતાર્થ અને બંને
સર્વપરિમાણકો યોગ્ય રીતે બંધ થાય.}

\begin{explain}
આ રીતથી ઉપગણો પર સીધું પરિમાણન કરવાને બદલે તેમના સંકેતો પર
પરિમાણન કરીને શક્તિસમૂહ વિશેનાં વિધાનો વ્યક્ત કરી શકીએ છીએ.
દાખલા તરીકે, કૅન્ટૉરનો પ્રમેય હવે એ રીતે કહી શકાય કે~$X$ના
શક્તિસમૂહને સંકેતિત કરતા કોઈપણ સંબંધના સંકેતોના ગણમાંથી~$X$
સુધી કોઈ !!{injective} વિધેય નથી.
\end{explain}

\begin{prop}
નીચેનું વાક્ય
\begin{multline*}
  \lforall[X][\lforall[Y][\lforall[R][(\fn{Pow}(Y, R, X) \lif \\
      \lnot \lexists[u][(\lforall[x][\lforall[y][(\eq[u(x)][u(y)] \lif
            \eq[x][y])]] \land {}\\
        \lforall[x][(Y(x) \lif X(u(x)))])])]]]
\end{multline*}
પ્રમાણભૂત છે.
\end{prop}

\begin{explain}
!!a{denumerable} ગણનો શક્તિસમૂહ !!{nonenumerable} હોય છે;
તેથી તેનું કાર્ડિનલ કદ કોઈપણ !!{denumerable} ગણના કદ~$\aleph_0$
કરતાં મોટું છે. $\Pow{\Nat}$ના કદને ``સાતત્યકનું કદ'' કહે છે,
કારણ કે તે વાસ્તવિક સંખ્યારેખા પરનાં બિંદુઓ, એટલે~$\Real$, જેટલું
છે. જો !!{domain}માં !!a{denumerable} ગણનો શક્તિસમૂહ સંકેતિત
કરવા જેટલા ઘટકો હોય, તો $\aleph_0$-કદના ગણ~$X$નો શક્તિસમૂહ
સંકેતિત કરતો કોઈપણ ગણ~$Y$ તેના જેટલા કદનો હોય એમ કહીને કોઈ ગણ
સાતત્યકના કદનો છે તે વ્યક્ત કરી શકીએ છીએ. (વ્યાપકક્ષેત્ર એટલું
મોટું ન હોય, એટલે તેમાં~$\Real$ સાથે સામ્ય કોઈ ઉપગણ ન હોય, તો
$\Pow{X}$ને સંકેતિત કરતો સંબંધ પણ હોઈ શકતો નથી.)
\end{explain}

\begin{prop}
જો $\cardle{\Real}{\Domain{M}}$ હોય, તો નીચેનું !!{formula}
\begin{multline*}
\fn{Cont}(Y) \ident
\lexists[X][\lexists[R][((\fn{Aleph}_0(X) \land
      \fn{Pow}(Y, R, X)) \land \\ 
      \lforall[x][\lforall[y][((Y(x) \land {}
      Y(y) \land \lforall[z][R(x,z) \liff R(y,z)]) \lif \eq[x][y])]])]]
\end{multline*}
$\cardeq{s(Y)}{\Real}$ વ્યક્ત કરે છે.
\end{prop}

\begin{proof}
$\fn{Pow}(Y, R, X)$ વ્યક્ત કરે છે કે~$s(Y)$ એ~$s(R)$ વડે
$s(X)$ના શક્તિસમૂહને સંકેતિત કરે છે; અને
$\fn{Aleph}_0(X)$ કહે છે કે તે ગણનીય અનંત છે. તેથી~$s(Y)$
ઓછામાં ઓછો સાતત્યકના કદનો છે, જોકે તે મોટો પણ હોઈ શકે છે
(જો~$s(Y)$ના ઘણા ઘટકો~$X$ના એક જ ઉપગણને સંકેતિત કરે).
છેલ્લો સંયોજક એ શક્યતા દૂર કરે છે: તે~$s(Y)$ના ઘટકો અને~$s(X)$ના
ઉપગણો વચ્ચેનો~$s(R)$ દ્વારા મળતો સંબંધ !!{injective} હોય એમ માગે છે.
\sourcecorrection{OLSOL-019}{સ્થિર અંગ્રેજી મૂળની સાબિતીમાં છેલ્લે
$s(Z)$ આવે છે, પરંતુ અહીં પરિમાણિત ગણ~$X$ છે; તેથી અહીં~$s(X)$
લખ્યું છે. મૂળનું ``be'' પણ વ્યાકરણભૂલ છે.}
\end{proof}

\begin{prop}
$\cardeq{\Domain{M}}{\Real}$ ત્યારે અને માત્ર ત્યારે જ જ્યારે
\begin{multline*}
  \Sat{M}{\lexists[Y][(\fn{Cont}(Y) \land \\
          \lexists[u][(\lforall[x][\lforall[y][(\eq[u(x)][u(y)] \lif \eq[x][y])]] \land {} \\
\lforall[y][(Y(y) \liff \lexists[x][\eq[y][u(x)]])])])]}.
\end{multline*}
\end{prop}
\sourcecorrection{OLSOL-020}{સ્થિર અંગ્રેજી મૂળનું અહીંનું વાક્ય
માત્ર~$Y$ શક્તિસમૂહને સંકેતિત કરે છે અને~$Y$ એ~$u$ની છબીનો
ઉપગણ છે એમ માગે છે; ઓળખ-વિધેયથી આ બીજી શરત સહેલાઈથી મળે છે,
અટલે વાક્ય મોટા વ્યાપકક્ષેત્રોમાં પણ સાચું થઈ શકે છે. વળી~$Y$
પર સંકેતીકરણ એક-એક છે તેવી શરત નથી. અહીં~$\fn{Cont}(Y)$થી
$Y$નું કદ સાતત્યક જેટલું નિશ્ચિત કર્યું છે અને~$u$નું મૂલ્યસમૂહ
બરાબર~$Y$ રાખ્યું છે, જેથી આખા વ્યાપકક્ષેત્રનું કદ પણ એટલું જ થાય.}

\begin{explain}
સાતત્યક પરિકલ્પના એવું વિધાન છે કે સાતત્યકનું કદ પહેલું
!!{nonenumerable} કાર્ડિનલ કદ છે, એટલે~$\Pow{\Nat}$નું
કદ~$\aleph_1$ છે.
\end{explain}

\begin{prop}
સાતત્યક પરિકલ્પના સાચી છે ત્યારે અને માત્ર ત્યારે જ જ્યારે
\[\fn{CH} \ident
\lforall[X][(\fn{Aleph}_1(X) \liff \fn{Cont}(X))]\]
પ્રમાણભૂત હોય.
\end{prop}

$\lnot \fn{CH}$ પ્રમાણભૂત હોય ત્યારે અને માત્ર ત્યારે જ સાતત્યક
પરિકલ્પના ખોટી છે, એવું માનવું સાચું નથી. !!a{enumerable}
વ્યાપકક્ષેત્રમાં~$\aleph_1$-કદના ઉપગણો પણ નથી અને સાતત્યકના
કદના ઉપગણો પણ નથી; તેથી !!a{enumerable} વ્યાપકક્ષેત્રમાં
$\fn{CH}$ હંમેશાં સાચું છે.
પરંતુ સાતત્યક પરિકલ્પના ખોટી છે ત્યારે અને માત્ર ત્યારે જ પ્રમાણભૂત
હોય એવું બીજું વાક્ય આપી શકીએ છીએ:

\begin{prop}
સાતત્યક પરિકલ્પના ખોટી છે ત્યારે અને માત્ર ત્યારે જ
\[\fn{NCH} \ident
\lforall[X][(\fn{Cont}(X) \lif \lexists[Y][(Y \subseteq X \land \lnot
    \fn{Count}(Y) \land \lnot \cardeq{X}{Y})])]\]
પ્રમાણભૂત હોય.
\end{prop}

\end{document}
